% Project brief for May 2026, Week 4.
% Date range: Mon 2026-05-25 to Sun 2026-05-31.
% Focus: RCA-Code model, public datasets, demo, and release.

\begin{frame}
  \titlepage
  \vspace{-1.0em}
  \begin{center}
    \small Project brief: May 2026, Week 4 \\
    \scriptsize Date range: Mon 2026-05-25 to Sun 2026-05-31
  \end{center}
\end{frame}

\begin{frame}{Background}
  \begin{itemize}
    \item RCA-Code targets debugging and trace-based root-cause analysis.
    \item Inputs can include issue text, stack traces, failing tests, code context, logs, patches, and tool observations.
    \item A direct long-context prompt is expensive and brittle for these evidence-heavy tasks.
    \item The project goal is a public model/demo, not just an internal experiment.
  \end{itemize}
\end{frame}

\begin{frame}{Project Significance}
  \begin{itemize}
    \item \textbf{Community impact}: release a usable RCA-Code model built on public data.
    \item \textbf{Industrial relevance}: debugging workflows already look like RCA over long traces and tool outputs.
    \item \textbf{Safety}: freeze the base model and train wrapper memory to avoid damaging original capabilities.
    \item \textbf{Demo value}: compare full context, summary, retrieval, and learned memory on the same bug/RCA case.
  \end{itemize}
\end{frame}

\begin{frame}{Project Scope and Boundary}
  \begin{itemize}
    \item Base: frozen 7B/8B coder or reasoning model.
    \item Trainable module: wrapper memory.
    \item Public coding/debug data: SWE-bench, BugsInPy, Defects4J, QuixBugs.
    \item RCA data: Nezha, OpenRCA-500, RCAEval, lincyaw/rca.
    \item RCA-demo initialization showed SFT can learn RCA signals, but direct SFT is not robust enough as the final RCA-Code route.
  \end{itemize}
\end{frame}

\begin{frame}{Benchmark Evidence from RCA-Demo}
  \begin{itemize}
    \item \textbf{Done / verified}: matrix evaluated 12 checkpoints across 4 RCA datasets and 5 general benchmarks.
    \item \textbf{SFT can help}: `q25_mix` reached `recommendation_f1 = 0.733` on RCAEval, about 18x over `q25_base`.
    \item \textbf{Curriculum has signal}: after fixing evaluation, `q25_curr_4_lincyaw` showed real service-identification gains on lincyaw.
    \item \textbf{But SFT is bounded}: q3 cumulative checkpoints showed format fragility; anti-forget v2 failed its retention gate.
  \end{itemize}
\end{frame}

\begin{frame}{Benchmark Takeaway for Project}
  \begin{itemize}
    \item Standard SFT is a useful initialization experiment, not the final RCA-Code architecture.
    \item We should not keep pushing base-model SFT as the only solution.
    \item Project direction: freeze the base model and move adaptation into wrapper memory.
    \item Immediate delivery remains a demo: compare full context, summary, retrieval, and wrapper memory on one public RCA/debug case.
  \end{itemize}
\end{frame}

\begin{frame}{First Delivery: Demo}
  \begin{enumerate}
    \item Pick one public RCA/debug case with long evidence.
    \item Compare full context, summary, retrieval, and wrapper memory.
    \item Show answer quality, evidence grounding, and context reduction.
    \item Use the demo to decide whether to scale to full RCA-Code benchmark runs.
    \item After the demo works, package model card, technical report, benchmark manifest, and release plan.
  \end{enumerate}
\end{frame}

\begin{frame}{Project Risks and Controls}
  \begin{itemize}
    \item Risk: learned memory compresses style but loses evidence grounding. Control: evidence metrics and qualitative failure cases.
    \item Risk: direct SFT hurts base capability. Control: frozen base in v0.
    \item Risk: private data leaks into release. Control: public-only release boundary; private DTS only internal.
    \item Risk: project becomes only a benchmark report. Control: ship model + demo + technical report.
  \end{itemize}
\end{frame}
